\section{Results}
\label{sec:results}

\subsection{Audit, size and representation}

\Cref{tab:audit} summarises the audit. The loader returns 2039 graphs from 2050 source records, the 11 missing entries being exactly those RDKit cannot parse, so no invalid or empty graph survives and no label is non-finite. No two raw SMILES strings in the file are identical, so string comparison reports zero duplicates, but canonicalisation reveals 60 molecules across 124 records, fifty repeating with a consistent label and ten carrying both (\cref{tab:conflicts}). \Cref{fig:molecules} shows six molecules with their SMILES strings and labels, and \cref{fig:distributions} the size distributions, which are unimodal and right-skewed around a median of 23 atoms and 25 bonds, with a tail reaching 132 atoms and bond count tracking atom count closely.

\begin{table}[htb]
    \centering
    \caption{Audit and molecule size over the 2039 loaded graphs. Bonds are counted once, not once per direction.}
    \label{tab:audit}
    \footnotesize
    \setlength{\tabcolsep}{3pt}
    \begin{tabular}{@{}lr@{\hspace{0.5em}}lrr@{}}
        \toprule
        Audit & Value & Size & Atoms & Bonds \\
        \midrule
        Loaded / in file      & 2039 / 2050 & Mean   & 24.06 & 25.95 \\
        Unparsable, dropped   & 11         & Std    & 10.59 & 11.72 \\
        Non-finite labels     & 0          & Min    & 2     & 1     \\
        Class 0 / class 1     & 479 / 1560 & Q1     & 18    & 19    \\
        Unique canonical      & 1975       & Median & 23    & 25    \\
        Duplicated molecules  & 60         & Q3     & 28.5  & 31    \\
        \quad consistent labels & 50       & Max    & 132   & 145   \\
        \quad conflicting labels & 10      &        &       &       \\
        \bottomrule
    \end{tabular}
\end{table}
Loading, encoding and batching were verified against the raw structures: the attribute matrices hold integer codes and not values, the COO list stores every bond twice, the batch index offsets are correct, and each batch carries one label per graph rather than per atom. Only a fraction of each vocabulary occurs --- 13 of 119 element slots, 4 of 12 formal charges, 4 of 8 hybridisation states and 4 of 22 bond types --- and encoding by \cref{eq:atom-encoder} with $d=64$ uses \num{11328} atom-embedding and \num{1920} bond-embedding parameters, where one-hot encoding would present 177 atom and 30 bond columns. Shapes, decoded vocabularies, the graph-versus-depiction comparison and the integrity checks are in \ref{sec:representation} and \ref{sec:integrity}.

\subsection{Splits}
\label{sec:res-splits}

Deduplication leaves \textbf{1965 unique molecules} at a positive rate of 0.7634 over 1020 scaffolds, the largest groups being benzene (128 molecules), the empty acyclic scaffold (95) and a steroid framework (74).

The label is not distributed uniformly through the file: in blocks of 250 records the positive rate is 0.656, 0.524, 0.532 and 0.428, then exactly 1.000 from row 1000 onwards, so the file is block-sorted by label in its second half. A greedy scaffold split breaking ties by file position therefore draws its singleton scaffolds from that all-positive tail and produces a single-class test set; shuffling equal-sized groups removes the artifact.

\Cref{tab:splits} reports both splits. Sizes and positive rates match the whole set to within half a percentage point, and both cover all 1965 molecules with disjoint indices. They differ in one number: the random split's subsets \textbf{share 238 scaffolds}, the scaffold split none.

\begin{table}[htb]
    \centering
    \caption{Composition of both splits over the 1965 unique molecules; the whole-set positive rate is 0.7634. Neither split shares a canonical SMILES between subsets; the random split shares \textbf{238 scaffolds}, the scaffold split none.}
    \label{tab:splits}
    \footnotesize
    \begin{tabular}{@{}llrrr@{}}
        \toprule
        Split & Subset & $n$ & Pos.\ rate & Scaffolds \\
        \midrule
        Random   & train & 1375 & 0.7636 & 773 \\
                 & val   &  295 & 0.7627 & 225 \\
                 & test  &  295 & 0.7627 & 226 \\
        \midrule
        Scaffold & train & 1375 & 0.7665 & 430 \\
                 & val   &  294 & 0.7585 & 294 \\
                 & test  &  296 & 0.7534 & 296 \\
        \bottomrule
    \end{tabular}
\end{table}
\subsection{Training}

Validation ROC-AUC for the baseline peaks at $C = 0.1$ on both splits (0.9024 random, 0.8875 scaffold) and falls monotonically to 0.8863 and 0.8576 at $C = 100$, so the sparse binary features need strong regularisation; the frozen thresholds differ, 0.534 and 0.650, because the validation probability distributions differ once validation molecules share no scaffold with training.

All six network runs early-stopped after 54--100 epochs; \cref{fig:curves} shows the curves and \cref{tab:seeds} the per-run outcome. Training loss falls smoothly to 0.08--0.12 while validation loss bottoms out near epoch 5--10 at about 0.25 and then climbs to 0.30--0.40, yet validation ROC-AUC does not degrade as that loss rises, which is why the restored checkpoints lie well past the validation-loss minimum. Seed variation is small --- $0.9154 \pm 0.0022$ and $0.8800 \pm 0.0025$ validation ROC-AUC --- while changing the split strategy moves the mean by 0.0354, about \textbf{seven times} the largest seed spread, with no overlap between the two groups: the worst random seed (0.9130) still beats the best scaffold seed (0.8822). Best epoch, by contrast, varies from 28 to 74.

\subsection{Test performance and errors}

\Cref{tab:results} reports the single evaluation of each sealed test subset, \cref{fig:roc-pr} the curves and \cref{fig:confusion} the confusion matrices. The network has the higher ROC-AUC and PR-AUC on both splits and the lower accuracy and $F_1$ on both, with thresholds spread widely across seeds; both models lose performance when scaffolds are held out, more so on the thresholded metrics (network $F_1$ 0.904 to 0.846).

\begin{table}[htb]
    \centering
    \caption{Test-set performance, one evaluation per sealed test subset. Network columns are mean $\pm$ standard deviation over three seeds; thresholds were frozen on validation data before any test molecule was scored. Better of the two models per split and metric in bold. A majority-class predictor scores ROC-AUC 0.500 at accuracy 0.763 and $F_1$ 0.865 (random), 0.753 and 0.859 (scaffold).}
    \label{tab:results}
    \scriptsize
    \setlength{\tabcolsep}{3pt}
    \begin{tabular}{@{}lcc@{\hspace{0.9em}}cc@{}}
        \toprule
        & \multicolumn{2}{c}{Random split} & \multicolumn{2}{c}{Scaffold split} \\
        \cmidrule(r){2-3} \cmidrule(l){4-5}
        & GINE & FP+LR & GINE & FP+LR \\
        \midrule
        Threshold   & $0.424 \pm 0.213$          & $0.534$          & $0.748 \pm 0.248$          & $0.650$          \\
        ROC-AUC     & $\mathbf{0.912 \pm 0.008}$ & $0.892$          & $\mathbf{0.885 \pm 0.025}$ & $0.876$          \\
        PR-AUC      & $\mathbf{0.967 \pm 0.005}$ & $0.955$          & $\mathbf{0.955 \pm 0.013}$ & $0.949$          \\
        Accuracy    & $0.859 \pm 0.029$          & $\mathbf{0.875}$ & $0.786 \pm 0.014$          & $\mathbf{0.821}$ \\
        Sensitivity & $0.881 \pm 0.052$          & $\mathbf{0.898}$ & $0.779 \pm 0.034$          & $\mathbf{0.834}$ \\
        Specificity & $0.786 \pm 0.052$          & $\mathbf{0.800}$ & $\mathbf{0.808 \pm 0.060}$ & $0.781$          \\
        $F_1$       & $0.904 \pm 0.022$          & $\mathbf{0.916}$ & $0.846 \pm 0.014$          & $\mathbf{0.875}$ \\
        \bottomrule
    \end{tabular}
\end{table}
The representative scaffold-split run (seed 2, frozen threshold 0.832) makes 68 errors on 296 test molecules, 12 false positives and 56 false negatives, the imbalance following from the high threshold. The median nearest-neighbour Tanimoto from the scaffold test set to the training set is \textbf{0.40}, well below the 0.7 convention, and the four most confident errors (\cref{tab:errors} and \cref{fig:errors}) all sit at or below it, between 0.20 and 0.38, three of them inheriting their nearest neighbour's label.
